\documentclass[aspectratio=169]{beamer}

\usepackage{tikz}

\usetheme{default}
\setbeamertemplate{navigation symbols}{}
\setbeamercovered{transparent=30}

\definecolor{pulpblue}{HTML}{347DBE}
\definecolor{pulpgreen}{HTML}{04AE19}
\definecolor{pulporange}{HTML}{FF9507}
\definecolor{pulpcharcoal}{HTML}{383838}
\definecolor{pulplightgreen}{HTML}{F4FAF2}

\setbeamercolor{background canvas}{bg=pulplightgreen}
\setbeamercolor{normal text}{fg=pulpcharcoal,bg=pulplightgreen}
\setbeamercolor{structure}{fg=pulpblue}
\setbeamercolor{title}{fg=pulpblue}
\setbeamercolor{subtitle}{fg=pulpgreen}
\setbeamercolor{author}{fg=pulpcharcoal}
\setbeamercolor{date}{fg=pulpcharcoal}
\setbeamercolor{frametitle}{fg=pulpblue,bg=pulplightgreen}
\setbeamercolor{itemize item}{fg=pulpgreen}
\setbeamercolor{itemize subitem}{fg=pulpblue}
\setbeamercolor{enumerate item}{fg=pulpblue}
\setbeamercolor{alerted text}{fg=pulporange}

\setbeamercolor{breadcrumb}{fg=black!65,bg=pulplightgreen}
\setbeamerfont{breadcrumb}{size=\scriptsize}
\setbeamertemplate{footline}{%
  \leavevmode%
  \hbox{%
    \begin{beamercolorbox}[
      wd=.82\paperwidth,
      ht=2.5ex,
      dp=3ex,
      leftskip=1em
    ]{breadcrumb}%
      \usebeamerfont{breadcrumb}%
      \ifnum\value{framenumber}=1
        {\small\insertshortauthor{} - 2026}
      \else\ifnum\value{section}>0
        \insertsectionhead%
        \ifnum\value{subsection}>0
          \enspace\textcolor{black!35}{/}\enspace
          \textbf{\insertsubsectionhead}%
        \fi
      \fi\fi
    \end{beamercolorbox}%
    \begin{beamercolorbox}[
      wd=.18\paperwidth,
      ht=2.5ex,
      dp=3ex,
      rightskip=1em
    ]{breadcrumb}%
      \usebeamerfont{breadcrumb}%
      \hfill\insertframenumber/\inserttotalframenumber
    \end{beamercolorbox}%
  }%
}

\title{Welcome Pulp Governance}
\subtitle{Past, Present and Future}
\author{Pedro Brochado}
\date{November 9, 2026}

\begin{document}

\begin{frame}
  \vspace*{\fill}
  \begin{center}
    \includegraphics[height=.18\paperheight]{_build/pulp-logo.pdf}\par
    \vspace{.1\paperheight}
    {\usebeamercolor[fg]{title}\fontsize{18}{18}\selectfont\bfseries
      \MakeUppercase{\inserttitle}\par}
    \vspace{.2em}
    {\usebeamercolor[fg]{subtitle}\Large\insertsubtitle\par}
  \end{center}
  \vspace*{\fill}
\end{frame}

\begin{frame}{Abstract}
  Summarize the governance question, the proposed approach, and the intended
  outcome.
\end{frame}

\section{Revision}

\begin{frame}{References reviewed}
  \begin{itemize}
    \item \textbf{Red Hat's CRA position}: industry context on lifecycle security.
    \item \textbf{Stewardship guidance}: Red Hat's stated role and practices for Pulp.
    \item \textbf{Community health audit (2018)}: historical governance and process context.
    \item \textbf{The Open Source Way}: governance roles, policies, and community evolution.
    \item \textbf{Pulp governance repository}: current project policies and processes.
  \end{itemize}
\end{frame}

\begin{frame}{Current context}
  Document existing governance arrangements. Cite primary sources for current
  roles, processes, and policies. Use the CRA materials for industry context
  and the 2018 audit for historical context.
\end{frame}

\section{Goal}

\begin{frame}{Goal}
  State the governance question, intended audience, and desired outcome.
\end{frame}

\section{What is governance?}

\begin{frame}{What is governance?}
  Define governance for the Pulp project and identify the roles, policies, and
  procedures that shape how the community works.
\end{frame}

\section{Past, Present and Future}

\subsection{Past}

\begin{frame}{Past}
  \begin{itemize}[<+->]
    \item Project health audits
    \item Community leadership
    \item Something else
  \end{itemize}
\end{frame}

\subsection{Present}

\begin{frame}{Present}
  \begin{itemize}
    \item Pulp governance repository
    \item Cyber Resilience Act (CRA)
  \end{itemize}
\end{frame}

\subsection{Future}

\begin{frame}{Future}
  \begin{itemize}
    \item Automated audits
    \item \texttt{MAINTAINERS.md} (CNCF)
  \end{itemize}
\end{frame}

\section{Open questions}

\begin{frame}{Open questions}
  Collect unresolved questions, potential risks, and decisions needed from the
  intended audience.
\end{frame}

\section{References}

\begin{frame}{References}
  \footnotesize
  \begin{itemize}
    \item \href{https://www.redhat.com/en/red-hat-eu-cyber-resilience-act}{Red Hat and the EU Cyber Resilience Act (CRA)}
    \item \href{https://access.redhat.com/security/eu-cyber-resilience-act-stewardship-guidelines}{EU Cyber Resilience Act Stewardship Guidelines}
    \item \href{https://pulpproject.org/blog/2018/09/17/pulp-community-health-audit-report/}{Pulp Community Health Audit Report (2018)}
    \item \href{https://guidebook.theopensourceway.org/growing-contributors/project-and-community-governance}{The Open Source Way: Project and Community Governance}
    \item \href{https://github.com/pulp/governance/}{Pulp governance repository}
  \end{itemize}
\end{frame}

\begin{frame}
  \begin{tikzpicture}[remember picture,overlay]
    \node[anchor=south,inner sep=0,yshift=4ex]
      at (current page.south)
      {\begin{tabular}{c}
        \includegraphics[height=.42\paperheight]{_build/pulp-mascot.pdf}\\[-.1em]
        {\small\textit{``Thanks!''}}
      \end{tabular}};
  \end{tikzpicture}
\end{frame}

\end{document}
